\subsection{Full column orbits and the three-by-three branch}

\begin{lemma}[Projective identifications of full column orbits]
\label{C:full-column-orbit-kernel}
Let $M$ have minimal band $I\times J$, and let $S$ be a finite
right act.  Suppose every core orbit in $SK$ has exactly $|J|$
elements.  If there are $o$ such orbits, the constant-action
quotient satisfies
\[
 0\longrightarrow(e\mathbb ZM)^{o-1}
 \longrightarrow\widetilde{\mathbb ZS}
 \longrightarrow\widetilde{\mathbb Z\bar S}
 \longrightarrow0,
\]
for a minimal idempotent $e$.  In particular the quotient
act-cone map is an integral homology isomorphism in every
degree at least three.
\end{lemma}
\begin{proof}
Label each full core orbit $O_a$ by its columns,
$s_{a,j}$, $j\in J$.  Each orbit is $M$-invariant and the
action on every orbit is the same column map:
$s_{a,j}m=s_{a,\tau_m(j)}$.
The relations of the constant-action quotient identify
$s_{a,j}$ with $s_{0,j}$ for fixed $j$, and make no other
identifications.  Its kernel therefore has integral basis
\[
 s_{a,j}-s_{0,j}\qquad(1\leq a<o,\ j\in J).
\]
For each $a$ these vectors form the right ideal module
$e\mathbb ZM$, proving the exact sequence and projectivity.
The act-cone Tor formula and its long exact sequence give the
claimed homology comparison, exactly as in
Lemma~\ref{C:two-column-projective-orbit-kernel}.
\end{proof}

\begin{theorem}[Three-by-three act cones below twenty-two arrows]
\label{C:three-by-three-act-cones21}
No act cone $\mathcal C(M,S)$ with minimal band $3\times3$
and $|M|+|S|+1\leq21$ realizes $\mathbb CP^2$.
\end{theorem}
\begin{proof}
Suppose otherwise and write $q=|M|-9$ for the number of
outside elements.  At most three states are excluded by the
earlier small-state results, so $q\leq7$ and $|S|\leq11-q$.
If $q\leq2$, Corollary~\ref{C:eleven-act-cones21} already
applies.  We may therefore assume $3\leq q\leq7$.

Let $c_j\in\bar S$ be the constant value of an old minimal
element in column $j$, and put $T=\{c_j:j\in J\}$.
If $|T|=1$, every old actor fixes this one state $s_0$.
For a balanced form $L$ on the compressed band, its general
affine representative is $u_L(f,g)=b^L_fa_g$, with
$b^L_m=L([s_0m]-[s_0],-)=0$ on every old actor.
This formula and its honest Karoubi transfer are the initial
construction in Theorem~\ref{C:act-common-fixed-row-cups};
that construction does not restrict the old column count.
The transferred two-cocycles thus
vanish on all pairs of old actors, making every degree-two
cup product zero on the act category.  This is impossible.

If $|T|=2$, identify old band columns with the same $c_j$,
while retaining all outside labels.  Equivariance
$c_jm=c_{\tau_m(j)}$ makes this a monoid quotient $M'$ with
minimal band $3\times2$ and the same $q$ outside elements.
The constant act $\bar S$ factors through $M'$.
The compression monoids for $(M,\bar S)$ and $(M',\bar S)$
are literally identical: their outside labels and products
agree, and their new band is $I\times\bar S$ with the same
actions and the same images of old band products.
Consequently the two act cones are homotopy equivalent through
that common monoid.  Since $|M'|=6+q\leq13$,
Theorem~\ref{C:three-row-fourteen-act-cones} makes every
rational degree-two product of this compressed cone zero,
without any finiteness assumption.  The degree-two quotient
surjection contradicts the original projective-plane square.

It remains that $|T|=3$.  Every core orbit has all three
columns: an equality between two column points in any core
orbit would identify the corresponding $c_j$ in $\bar S$.
Apply Lemma~\ref{C:full-column-orbit-kernel} and compress to
$N$.  Its rational homology is bounded and
$H_4(BN;\mathbb Q)=\mathbb Q$, because both are true of the
original cone and the homology comparison holds in degrees
at least three.  The contracted algebra still has dimension
$q$.  If $q\leq4$, the band Tor formula and
Theorem~\ref{B:contracted-four-Tor} give an immediate
contradiction.

Thus $5\leq q\leq7$ and $|S|\leq6$.  The act is connected:
the first homology sequence of the act-cone pushout identifies
$H_1$ with the reduced degree-zero homology of its action
space, since $H_1(BM;\mathbb Q)=0$.
Two full core orbits would already use six states, and each
is individually $M$-invariant.  At this state bound they would
exhaust $S$ and disconnect the act.  There is therefore only
one core orbit.  Every band actor is then constant on all of
$S$: its image lies in that orbit, and idempotence shows it
equals the constant value of its restriction to the orbit.
Compression is now an actual homotopy equivalence from the
original cone to a monoid with
\[
 |N|=q+3|S|\leq33-2q.
\]
For $q=6,7$ this is at most $21$, contrary to the exact
monoid minimum $23$.  For $q=5$, it is at most $23$;
the only case not already excluded by the monoid minimum is
$|S|=6$.  Its compressed monoid has minimal band $3\times6$
and five outside elements, excluded by
Lemma~\ref{C:three-by-six-five-outside-CP2}.
This exhausts the cases.
\end{proof}
